% META: {"id": "TNSI-BDD-EVAL-A", "chapitre": "TNSI-BASES-DE-DONNEES", "type_objet": "evaluation", "capacites": ["T-BDD-01A", "T-BDD-03B", "T-BDD-03C", "T-BDD-03F"], "mode_creation": "ex_nihilo", "sources_inspiration": [], "duree_min": 45, "status": "approved"}
\section*{Évaluation A --- Bases de données}
\textit{Durée : 45 minutes. Ordinateur avec un SGBD (SQLite ou équivalent) autorisé.}

\baremeIndicatif{Ex. 1 : 8 pts (modèle relationnel) ; Ex. 2 : 12 pts (SELECT, WHERE, INSERT)}

\bigskip

\textbf{Exercice 1} --- \textit{Modélisation} \hfill (8 points)

\begin{enumerate}
  \item (4 pts) Définir, en une phrase chacun, les termes \textit{relation}, \textit{attribut} et \textit{clé primaire}.
  \item (4 pts) Donner un exemple de clé étrangère dans le schéma \texttt{Commande(id, id\_client, date)} / \texttt{Client(id, nom)}, et expliquer à quoi elle sert.
\end{enumerate}

\bigskip

\textbf{Exercice 2} --- \textit{SELECT, WHERE, INSERT} \hfill (12 points)

On donne la table :
\begin{sql}
CREATE TABLE Produit(id INTEGER PRIMARY KEY, nom TEXT, prix REAL, categorie TEXT);
INSERT INTO Produit VALUES
  (1, 'Cahier', 1.50, 'Papeterie'),
  (2, 'Stylo', 0.80, 'Papeterie'),
  (3, 'Casque audio', 25.00, 'Electronique');
\end{sql}
\begin{enumerate}
  \item (4 pts) Écrire la requête affichant le nom des produits de catégorie \texttt{'Papeterie'}.
  \item (4 pts) Écrire la requête affichant le nom des produits dont le prix est strictement inférieur à $2.00$.
  \item (4 pts) Écrire la requête ajoutant un produit \texttt{'Gomme'}, catégorie \texttt{'Papeterie'}, prix $0.60$.
\end{enumerate}

% BEGIN-VERIFY
% import sqlite3
% conn = sqlite3.connect(":memory:")
% conn.execute("CREATE TABLE Produit(id INTEGER PRIMARY KEY, nom TEXT, prix REAL, categorie TEXT)")
% conn.executemany("INSERT INTO Produit VALUES (?,?,?,?)", [
%     (1, "Cahier", 1.50, "Papeterie"), (2, "Stylo", 0.80, "Papeterie"), (3, "Casque audio", 25.00, "Electronique"),
% ])
% conn.commit()
% cur = conn.execute("SELECT nom FROM Produit WHERE categorie = 'Papeterie'")
% assert set(r[0] for r in cur.fetchall()) == {"Cahier", "Stylo"}
% cur = conn.execute("SELECT nom FROM Produit WHERE prix < 2.00")
% assert set(r[0] for r in cur.fetchall()) == {"Cahier", "Stylo"}
% conn.execute("INSERT INTO Produit(nom, categorie, prix) VALUES ('Gomme', 'Papeterie', 0.60)")
% conn.commit()
% assert conn.execute("SELECT COUNT(*) FROM Produit").fetchone()[0] == 4
% END-VERIFY
